\documentclass[11pt,a4paper]{article}
\newcommand{\labonaamwaarde}{labo 3}
\usepackage{syllabus}

\begin{document}

\syllabustitel{labo 3 -- \emph{Dataverwerking I}}{Van bestand tot nette tabel}

% ---------------------------------------------------------------------------
% Gedeeld door de syllabi van labo 3, 4 en 5. Pas dit hier aan, niet in de labobestanden.
\section{Overzicht}

Tijdens labo 3, 4 en 5 ben jij de nieuwe data-analist van een
(fictieve) jeugdclub voor voetbal en atletiek. De hoofdtrainer bezorgt je
drie exportbestanden uit de systemen van de club. Die zijn bewust
rommelig, zoals echte exports: inconsistente labels, sprinttijden in
milliseconden, ontbrekende waarden, een atleet die twee keer voorkomt en
een kolom met een komma als decimaalteken.

Elke sessie bouwt verder op de vorige:

\begin{description}[leftmargin=1.6cm, style=nextline]
  \item[Labo 3] \textbf{Van bestand tot nette tabel.} Inladen, verkennen,
    filteren, kolommen toevoegen en opkuisen tot je de data kunt
    vertrouwen.
  \item[Labo 4] \textbf{Combineren, samenvatten, herschikken.} Groepen
    samenvatten, tabellen koppelen, wisselen tussen een brede en een lange
    vorm, en werken met een hartslagreeks per seconde.
  \item[Labo 5] \textbf{Vertel het verhaal met plotly.} De juiste grafiek
    kiezen voor een vraag, eerlijke grafieken maken met
    \code{plotly.express} en ze combineren tot een klein interactief
    dashboard.
\end{description}


% ---------------------------------------------------------------------------
\section{Leerdoelen}

Na labo 3 kun je:

\begin{enumerate}[label=\textbf{LD\arabic*}, leftmargin=1.4cm]
  \item een CSV-bestand inladen met \code{pandas} en de structuur, de
        datatypes en de ontbrekende waarden beschrijven;
  \item rijen en kolommen selecteren en filteren met logische
        voorwaarden, en het resultaat sorteren;
  \item nieuwe kolommen afleiden (leeftijd, snelheid, BMI) en tekst
        omzetten naar getallen en datums;
  \item een dataset opkuisen en elke keuze verantwoorden: inconsistente
        tekstlabels, verkeerde eenheden, ontbrekende waarden en dubbels.
\end{enumerate}

% ---------------------------------------------------------------------------
\section{Voorkennis}

Labo 1 en 2 (\emph{Introductie Python I en II}): variabelen en datatypes,
operatoren, \code{if}/\code{elif}/\code{else}, functies,
stringmethodes, lijsten, dictionaries, lussen en een CSV-bestand inlezen
met de module \code{csv}. Labo 2 eindigt met een samenvatting van
\code{athletes.csv} in gewone Python; in labo 3 zie je hoe \code{pandas}
hetzelfde doet in enkele regels.

% ---------------------------------------------------------------------------
% Gedeeld door de syllabi van labo 3, 4 en 5. Pas dit hier aan, niet in de labobestanden.
\section{Software en lesmateriaal}

Elk labo is een \textbf{marimo-notebook} in de map \code{labs/} van de
cursusrepository. Je kunt het op twee manieren openen:

\begin{itemize}
    \item \textbf{Online:} surf naar \url{https://marimo.app}, kies
        \emph{New $\rightarrow$ Open from URL\dots} en plak de raw-URL van
        het notebook op GitHub. De data wordt dan van GitHub gelezen, dus
        je hoeft niets te installeren.
    \item \textbf{Lokaal:} installeer de pakketten uit
        \code{requirements.txt} en open het notebook met marimo. De data
        wordt dan uit de map \code{data/} gelezen.

\end{itemize}

Drie dingen die je moet weten over marimo:

\begin{itemize}
  \item Notebooks zijn \textbf{reactief}: als je een cel aanpast, worden
        alle cellen die ervan afhangen opnieuw uitgevoerd, in de volgorde
        van hun afhankelijkheden.
  \item \textbf{Twee cellen kunnen niet dezelfde variabele definiëren.}
        Geef elke stap een eigen, duidelijke naam (\code{tests\_raw},
        \code{tests\_clean}, \dots) of zet een reeks stappen samen in één
        cel.
  \item Elke oefening bevat een placeholder \code{...}. Een verborgen cel
        eronder controleert je antwoord: die kleurt \textbf{groen} als het
        juist is, en geeft anders een algemene hint over wat nog niet
        klopt, zonder het antwoord te verklappen.
\end{itemize}


% ---------------------------------------------------------------------------
% Gedeeld door de syllabi van labo 3, 4 en 5. Pas dit hier aan, niet in de labobestanden.
\section{De dataset}

\noindent
\rowcolors{2}{tint}{white}
\begin{tabularx}{\textwidth}{@{}L{3.1cm}L{3.2cm}Y@{}}
\toprule
\textbf{Bestand} & \textbf{Omvang} & \textbf{Inhoud en gekende problemen} \\
\midrule
\code{athletes.csv}   & 43 rijen, 7 kolommen   & ID, naam, geboortedatum, geslacht, positie, lengte, gewicht. 13 schrijfwijzen voor 4 posities, gewicht als tekst met een decimale komma (\code{"57,5"}), één atleet twee keer ingevoerd met twee ID's, enkele ontbrekende waarden. \\
\code{tests.csv}      & 114 rijen, 8 kolommen  & Sprint over 10\,m en 30\,m, spronghoogte (countermovement jump) en Yo-Yo-niveau, op 3 testmomenten. Sommige sprinttijden staan in milliseconden; enkele ontbrekende waarden. \\
\code{hr\_session.csv} & 12\,000 rijen, 5 kolommen & Hartslag en snelheid per seconde van 5 atleten tijdens één training van 40 minuten. \\
\bottomrule
\end{tabularx}
\rowcolors{1}{}{}


% ---------------------------------------------------------------------------
\needspace{18\baselineskip}
\section{Planning}

% Planning van één labo. Gebruikt door de syllabus van dit labo en door de syllabus voor labo 3-5.
\textbf{Doel:} een dataset inladen en verkennen, rijen en kolommen
selecteren en filteren, nieuwe kolommen toevoegen en de data opkuisen.

\begin{schedule}
13:00--13:20 & Opwarmer: waarom niet gewoon Excel? & De ruwe tekst van een CSV-bestand bekijken en zien wat een tabelbibliotheek toevoegt. \\
13:20--14:00 & Inladen en verkennen & \code{read\_csv}, \code{.head()}, \code{.shape}, \code{.dtypes}, \code{.info()}, \code{.describe()}, \code{mo.ui.table}. \\
14:00--14:50 & Selecteren en filteren & Kolommen selecteren, \code{.loc} met logische voorwaarden, voorwaarden combineren, \code{sort\_values}. \\
14:50--15:05 & Pauze & \\
15:05--15:50 & Nieuwe kolommen & Tekst naar getal, leeftijd uit een geboortedatum, snelheid uit een sprinttijd, BMI. \\
15:50--17:00 & Opkuisen & Ontbrekende waarden (\code{dropna} of \code{fillna}), tekstlabels, eenheden, dubbels, één volledige opkuis-pipeline. \\
17:00--17:30 & Uitdaging & De 5 snelste atleten per positie. Exit ticket. \\
\end{schedule}

\textbf{Oefeningen:} (1) \code{tests.csv} inladen; (2) aanvallers jonger
dan 16; (3) verdachte sprinttijden; (4) BMI; (5) leeftijd in jaren;
(6) de posities opkuisen; (7) de eenheden van de sprinttijden
corrigeren; (8) de dubbele atleet verwijderen; \textbf{uitdaging:} de
5 snelsten per positie.


% ---------------------------------------------------------------------------
\section{Werkwijze}

\begin{itemize}
  \item Elk blok start met een korte demo in het notebook, gevolgd door
        oefeningen die je in het notebook zelf oplost.
  \item Gebruik de automatische controles als feedback, niet als doel:
        als een controle groen kleurt, zorg dan dat je kunt uitleggen
        \emph{waarom} je code werkt.
  \item Het labo eindigt met een exit ticket: één techniek die je opnieuw
        zult gebruiken en één ding dat nog niet duidelijk is. Breng de
        onduidelijke punten mee naar labo 4.
\end{itemize}

% ---------------------------------------------------------------------------
\section{Voorbereiding}

\begin{itemize}
  \item Controleer vóór de sessie dat je het notebook kunt openen,
        lokaal of op marimo.app.
  \item Herhaal lijsten, dictionaries en functies uit labo 2: pandas
        gebruikt ze voortdurend.
  \item Lees de cursustekst bij labo 3 naast het notebook.
\end{itemize}

\end{document}
